\documentclass{article}
\usepackage[T1]{fontenc}
\usepackage{lmodern}
\usepackage{graphicx}
\usepackage{amsmath,amssymb}
\usepackage[a4paper,margin=2cm]{geometry}
\usepackage[absolute,overlay]{textpos}
\setlength{\TPHorizModule}{1mm}
\setlength{\TPVertModule}{1mm}
\usepackage[indonesian]{babel}
\usepackage{hyperref}
\setlength{\emergencystretch}{2em}
% unit: Halaman 14

\begin{document}

% === Halaman 14 (llm) ===

\textbf{Contoh 1:} Bob membangkitkan kunci publik dan kunci privatnya. Alice akan mengenkripsi pesan dengan menggunakan kunci publik Bob.

\textbf{(a) Pembangkitan kunci} (Oleh Bob)

\textcolor{red}{syarat $g < p, g$ akar primitif dari 2357 dan $1 \le x \le p - 2$}

Misal $p = 2357, g = 2,$ dan $x = 1751$.

Hitung: $y = g^x \pmod p = 2^{1751} \pmod{2353} = 1185$

Hasil: Kunci publik: $(y = 1185, g = 2, p = 2357)$

Kunci privat: $(x = 1751, p = 2357)$.

Bob memberitahu kunci publik ini kepada Alice

\vspace{10mm}

\textbf{(b) Enkripsi} (Oleh Alice)

Misalkan pesan $m = 2035$ (nilai $m$ masih berada di dalam selang $[0, 2357 - 1])$.

Alice memilih bilangan acak $k = 1520$ (nilai $k$ berada di dalam selang $[0, 2357 - 1])$.

\end{document}
